% !TEX root = ../documentation.tex
\chapter{Building a Theme}\label{S:ThemeCommands}


	This part of the manual documents the commands used to build themes. They are intended for designers creating or modifying a theme, and for users who wish to adjust an existing theme in detail. By convention, all of these commands begin with \texttt{\textbackslash{}RpgTheme...}.

	\begin{center}
		\large \bfseries Authors of \rpgtex{} documents probably do not need any of the commands in this part, and it can safely be ignored.
	\end{center}
	This chapter describes how a theme is organised. The chapters which follow describe how a theme defines its fonts \RpgPage[p]{C:Fonts}, the design of its pages \RpgPage[p]{C:PageDesign}, the appearance of individual elements \RpgPage[p]{S:ElementStyles}, and its rule environments \RpgPage[p]{C:ThemeRuleEnvs}.

	\begin{RpgSidebar}{Code Setters}
		Many of the commands in this part of the manual do not perform any action themselves. Instead, they \emph{store} a piece of code, which \rpgtex{} runs later, when the corresponding element is created. For example, \cmd{RpgThemePartPage} stores the code used to draw a part page, which is run each time \cmd{part} is called.

		The stored code is often given arguments at the time it is run, such as the name of the part. These are accessed as \texttt{\#1}, \texttt{\#2} and so on within the stored code, even though the setter itself is called with only a single argument:
		\begin{Code}
			\cmd{RpgThemePartPage}\{The part is called \#1\}
		\end{Code}
		The entry for each setter lists the arguments which are available.
	\end{RpgSidebar}

	\section{Theme Files}\label{S:ThemeStructure}

		When \cmd{RpgSetTheme}\param{<name>} is called, \RpgPage[p]{Cmd:RpgSetTheme} \rpgtex{} attempts to input a file of the form:
		\begin{Code}
			<theme path> / <name> / <name>.cfg
		\end{Code}
		Where \texttt{<theme path>} is the value set by the theme path option \RpgPage[p]{S:ThemePath} or by \cmdref{RpgSetThemePath}. The theme ``\texttt{name}'' is therefore defined by a directory, \texttt{<theme path>/name}, containing the file \texttt{name.cfg}. This file (the `theme configuration file') may then execute any \LaTeX{} code.

		If the theme config file separates out the concerns across multiple files, it is important to \cmd*{input} them using the full file path to prevent name collisions. For example, the default theme's configuration file begins:
		\begin{Code}
			\cmd*{input}\{\cmd*{RpgPackagePath}/themes/default/colors.theme.tex\}

			\cmd*{input}\{\cmd*{RpgPackagePath}/themes/default/utility.theme.tex\}

			\ldots
		\end{Code}

		\subsection{Theme Code}\label{S:ThemeCode}

			Since themes are inserted entirely through \cmd*{input} commands, and may be introduced midway through a document, there are some restrictions and limitations on what should go inside a theme file or its subsidiary files:

			\begin{enumerate}
				\item \textbf{Packages.} Any package loading content (\cmd*{usepackage}, \cmd*{tcbuselibrary}) should be wrapped in a \cmdref{RpgIfPreamble} block, and appropriate guards put in place in case the packages cannot be loaded (because the theme was first activated in the document body).
				\item \textbf{Other preamble-only commands.} Commands such as \cmd*{geometry} may also only be used in the preamble. Where a body equivalent exists, \cmdref{RpgIfPreambleTF} can choose between the two (for example, \cmd*{geometry} in the preamble and \cmd*{newgeometry} in the body).
				\item \textbf{Repeated loading.} A theme may be loaded more than once in the same document, either because the author switches back to it, or because another theme builds upon it: the \texttt{dnd} and \texttt{scifi} themes both begin by calling \cmd{RpgSetTheme}\{default\}, and so the default theme is reloaded whenever either is activated. Every definition must therefore succeed when it is made a second time:
				      \begin{itemize}
					      \item Use \cmd*{DeclareDocumentCommand} (or \cmd*{cs_set:Npn}) rather than \cmd*{NewDocumentCommand} (or \cmd*{cs_new:Npn}), so that the definition is replaced rather than raising an error. This also allows a later theme to override a command defined by an earlier one.
					      \item Create variables only if they do not already exist (for example, with \cmd*{tl_if_exist:NF}), and use \cmd*{msg_set:nnn} rather than \cmd*{msg_new:nnn}.
				      \end{itemize}
				\item \textbf{Category codes.} A theme file is read with the category codes in force where \cmd{RpgSetTheme} is called. When it is called in the document body, \LaTeX{}3 syntax is off and \texttt{@} is not a letter. Each file should therefore begin with \cmd*{ExplSyntaxOn} if it uses \LaTeX{}3 syntax, and with \cmd*{makeatletter} if it uses internal commands such as \cmd*{@title}, rather than relying on the settings left by the file before it.
				\item \textbf{Timing.} A theme only affects material created after it is loaded. In particular, it cannot change the options passed to the underlying class (such as the font size or the number of columns), and it does not restyle a table of contents which has already been typeset.
			\end{enumerate}

		\subsection{Theme Resources}\label{S:ThemeResources}

			It is possible to ship non-\TeX{} resources, such as fonts and images, with a theme. By convention, these are kept in subdirectories of the theme's directory, such as \texttt{fonts/} and \texttt{img/}.

			\TeX{}, \texttt{graphicx} and \texttt{fontspec} search for these files relative to the document being compiled, not to the theme, and so each resource must be referred to by its full path. For example, the \texttt{dnd} theme sets its paper texture with:
			\begin{Code}
				\cmd*{RpgThemePaperImage}\{\cmd*{RpgPackagePath}/themes/dnd/img/paper\}
			\end{Code}

			Note that this uses \cmd{RpgPackagePath} because the theme is bundled with a package; for a general theme which is stored elsewhere, the designer must give the full path.

			Fonts are loaded by giving \texttt{fontspec} the location of the font files. For example, the \texttt{scifi} theme loads its body font with:
			\begin{Code}
				\cmd*{RpgIfFont}\{

				\quad\cmd*{providefontfamily}\{\cmd*{jura}\}\{Jura\}[

				\quad\quad Path = \cmd*{RpgPackagePath}/themes/scifi/fonts/,

				\quad\quad UprightFont = *-Regular,

				\quad\quad \ldots

				\quad]

				\}
			\end{Code}
			\cmdref{RpgThemeFont} may be used unguarded, since it does not load any fonts itself, but the font declarations should be placed within \cmdref{RpgIfFont}, so that fonts are only loaded when the \texttt{font} option is active.

	\section{Conditional Code}\label{S:ThemeConditional}

		Many parts of a theme are only meaningful when the corresponding package option is active, and some commands may only be used in the preamble, although a theme may be loaded in the document body. The following commands allow a theme to include code conditionally:

		\RpgFunction{RpgIfFontTF}[RpgIfFont]{\param{true code}\param{false code}}
		{
			Runs \texttt{true code} if the \texttt{font} option is active \RpgPage[p]{S:PackageOptions}, and \texttt{false code} otherwise. The alias \cmd{RpgIfFont}\param{code} takes a single argument, and is equivalent to \cmd{RpgIfFontTF}\param{code}\{\}.
		}
		\RpgFunction{RpgIfPaperTF}[RpgIfPaper]{\param{true code}\param{false code}}
		{
			Runs \texttt{true code} if the \texttt{background} option is active, and \texttt{false code} otherwise. The alias \cmd{RpgIfPaper}\param{code} takes a single argument, and is equivalent to \cmd{RpgIfPaperTF}\param{code}\{\}.
		}
		\RpgFunction{RpgIfPreambleTF}[RpgIfPreamble]{\param{true code}\param{false code}}
		{
			Runs \texttt{true code} if it is reached before \texttt{\textbackslash begin\{document\}}, and \texttt{false code} otherwise. The alias \cmd{RpgIfPreamble}\param{code} takes a single argument, and is equivalent to \cmd{RpgIfPreambleTF}\param{code}\{\}.

			This allows a body alternative to be given for preamble-only commands, for example:
			\begin{Code}
				\cmd*{RpgIfPreambleTF}\{\cmd*{geometry}\{...\}\}\{\cmd*{newgeometry}\{...\}\}
			\end{Code}
			The restrictions on what a theme may contain are described on \RpgPage{S:ThemeCode}.
		}
		\RpgFunction{RpgIfSectionsTF}[RpgIfSections]{\param{true code}\param{false code}}
		{
			Runs \texttt{true code} if the \texttt{sections} option is active, and \texttt{false code} otherwise. The alias \cmd{RpgIfSections}\param{code} takes a single argument, and is equivalent to \cmd{RpgIfSectionsTF}\param{code}\{\}.
		}
		\RpgFunction{RpgIfTitleTF}[RpgIfTitle]{\param{true code}\param{false code}}
		{
			Runs \texttt{true code} if the \texttt{title} option is active, and \texttt{false code} otherwise. The alias \cmd{RpgIfTitle}\param{code} takes a single argument, and is equivalent to \cmd{RpgIfTitleTF}\param{code}\{\}.
		}

